\documentclass[12pt]{article}
\usepackage[T1]{fontenc}
\usepackage[utf8]{inputenc}
\usepackage{lmodern}
\usepackage{textcomp}
\usepackage{enumitem}
\usepackage{geometry}
\geometry{margin=1in}
\begin{document}
\begin{center}
Answer Key - Economics - Graduate Level Assessment 20 \
\small (Answers are ordered exactly as in the question paper.)
\end{center}

\section*{Multiple Choice}
\begin{enumerate}[label=\arabic*.]
\item \textbf{Answer:} C
\\ \textit{Options:} A) Bill rate is 7.70\%, Coupon issue yield is 8.50\%; B) Bill rate is 8.50\%, Coupon issue yield is 8.24\%; C) Bill rate is 7.96\%, Coupon issue yield is 8.24\%; D) Bill rate is 8.24\%, Coupon issue yield is 8.50\%; E) Bill rate is 8.24\%, Coupon issue yield is 7.96\%
\\ \textit{Note:} The bill rate is calculated using the formula that considers the discount from the face value to the purchase price over the maturity period, resulting in 7.96\%, while the coupon-issue yield equivalent accounts for the annualized yield based on the purchase price and the discount, yielding 8.24\%.
\item \textbf{Answer:} A
\\ \textit{Options:} A) P = MC.; B) Perfect competition.; C) Efficiency.; D) Inelastic demand.; E) P =MR.
\\ \textit{Note:} In monopolistic competition, firms set prices above marginal cost (P \textgreater{} MC) due to product differentiation, but the answer "P = MC" indicates the condition for allocative efficiency, which is not typically achieved in this market structure, highlighting the distinction between theoretical efficiency and practical market behavior.
\item \textbf{Answer:} A
\\ \textit{Options:} A) A government regulation.; B) An increase in supply.; C) A government tax.; D) A government fine.; E) A decrease in demand.
\\ \textit{Note:} A government regulation is correct because it can incentivize or require individuals or firms to produce more of a good with positive externalities, thus aligning private benefits with social benefits and correcting the market failure.
\item \textbf{Answer:} C
\\ \textit{Options:} A) price = marginal revenue.; B) price = minimum average cost.; C) price = minimum average variable cost.; D) price = marginal cost.
\\ \textit{Note:} In the long run, a perfectly competitive firm operates where price equals the minimum average total cost, not the minimum average variable cost, as firms must cover all costs to remain in the market.
\item \textbf{Answer:} A
\\ \textit{Options:} A) Net exports and GDP both fall.; B) GDP falls and there is no change in net exports.; C) Net exports rise and there is no change in GDP.; D) Net exports and GDP both rise.; E) Net exports rise and GDP falls.
\\ \textit{Note:} The purchase of an entertainment system manufactured in China is considered an import, which decreases net exports and, consequently, reduces the gross domestic product (GDP) of the U.S. economy.
\end{enumerate}

\section*{True / False}
\begin{enumerate}[label=\arabic*.]
\setcounter{enumi}{5}
\item \textbf{Answer:} False
\\ \textit{Options:} A) True; B) False
\\ \textit{Note:} A trade deficit of $603 billion indicates that the United States imported $603 billion more in goods and services than it exported, reflecting a greater volume of imports compared to exports.
\item \textbf{Answer:} False
\\ \textit{Options:} A) True; B) False
\\ \textit{Note:} The statement is false because economic resources, or factors of production, traditionally include land, labor, capital, and entrepreneurship, while money, machines, and management are merely components or tools used to utilize these resources effectively.
\item \textbf{Answer:} False
\\ \textit{Options:} A) True; B) False
\\ \textit{Note:} The primary language spoken in a country does not directly influence export and import volumes, as these volumes are primarily determined by factors such as economic policies, trade agreements, resource availability, and market demand rather than linguistic factors.
\item \textbf{Answer:} False
\\ \textit{Options:} A) True; B) False
\\ \textit{Note:} The statement is false because a growth rate of 10 percent per year would actually result in the standard of living doubling in approximately 7.2 years, not 12, according to the Rule of 72.
\item \textbf{Answer:} True
\\ \textit{Options:} A) True; B) False
\\ \textit{Note:} Marginal Fixed Cost is zero because fixed costs do not change with the level of output, while Marginal Variable Cost is equal to Marginal Total Cost since total cost changes solely with variable costs at the margin.
\end{enumerate}

\section*{Long Form Response}
\begin{enumerate}[label=\arabic*.]
\setcounter{enumi}{10}
\item \textbf{Answer:} Positive externalities associated with the production or consumption of a good lead to market failure when the marginal social benefit (MSB) exceeds the marginal social cost (MSC). This discrepancy indicates that the benefits to society from the good are not fully reflected in its market price, resulting in underproduction or underconsumption relative to the socially optimal level. For instance, in the case of education, the societal gains from an educated populace, such as increased productivity and reduced crime rates, extend beyond the individual's private benefits. Consequently, without appropriate intervention or subsidies to align MSB with MSC, resources are misallocated, and overall welfare is not maximized, illustrating a fundamental market failure.
\\ \textit{Note:} correct because it accurately describes how positive externalities create a situation where the marginal social benefit exceeds the marginal social cost, leading to underproduction or underconsumption of goods, which ultimately results in market failure.
\item \textbf{Answer:} The imposition of an excise tax on the production of good X in a competitive market creates a shift in the burden of the tax between consumers and producers, heavily influenced by the elasticity of demand. When demand for good X is inelastic, consumers are less responsive to price changes, allowing producers to pass a larger portion of the tax onto consumers in the form of higher prices. Conversely, if demand is elastic, consumers can significantly reduce their quantity demanded in response to price increases, which forces producers to absorb more of the tax burden. Thus, the distribution of the tax burden ultimately reflects the relative elasticities of demand and supply, with inelastic demand leading to a greater consumer burden. This phenomenon underscores the importance of understanding elasticity when analyzing the economic effects of taxation in competitive markets.
\\ \textit{Note:} correct because it accurately describes how the burden of an excise tax is shared between consumers and producers, emphasizing that the degree of demand elasticity determines the extent to which producers can pass on the tax to consumers through price increases.
\end{enumerate}

\vfill
\noindent\textit{End of Answer Key}
\end{document}